%% ma: L11-B04-0011
%% lop: 11
%% chuong: 1
%% bai: 4
%% dang: 11.4.10
%% loai: TL
%% muc: VD
%% dap_an: "-"
%% nguon: "de-cuong-giua-ky-1-toan-11-nam-2026-2027-truong-thpt-hoang-van-thu-ha-noi-2-8.pdf (M08, Đề cương giữa kì 1 Toán 11 THPT Hoàng Văn Thụ, Hà Nội), Phần 4 Bài 4"
%% kien_thuc: "Bài 2 (công thức lượng giác) (biến đổi tổng thành tích ở ý d), Bài 4 (phương trình lượng giác cơ bản)"
%% trang_thai: chinh-thuc
%% ghi_chu: "Đã duyệt 29/10/2026."
\begin{ex}
Giải các phương trình sau:

a) $\tan\left(\dfrac{x}{2}+\dfrac{\pi}{6}\right)=\sqrt3$;\quad b) $\sin\left(3x+\dfrac{\pi}{4}\right)=\dfrac{\sqrt3}{2}$;

c) $\cos(2x-70^\circ)=-1$;\quad d) $\cos\left(x+\dfrac{\pi}{4}\right)+\cos\left(\dfrac{\pi}{4}-x\right)=0$.
\loigiai{a) $\dfrac x2+\dfrac\pi6=\dfrac\pi3+k\pi\Leftrightarrow x=\dfrac\pi3+k2\pi$ $(k\in\mathbb Z)$.

b) $3x+\dfrac\pi4=\dfrac\pi3+k2\pi$ hoặc $3x+\dfrac\pi4=\dfrac{2\pi}3+k2\pi\Leftrightarrow x=\dfrac\pi{36}+\dfrac{k2\pi}3$ hoặc $x=\dfrac{5\pi}{36}+\dfrac{k2\pi}3$ $(k\in\mathbb Z)$.

c) $2x-70^\circ=180^\circ+k360^\circ\Leftrightarrow x=125^\circ+k180^\circ$ $(k\in\mathbb Z)$.

d) $2\cos\dfrac\pi4\cos x=0\Leftrightarrow\cos x=0\Leftrightarrow x=\dfrac\pi2+k\pi$ $(k\in\mathbb Z)$.}
\end{ex}
